\documentclass{article}
\usepackage{graphicx} % Required for inserting images
\usepackage{amsmath}
\usepackage{amssymb}
\usepackage{amsthm}
\newtheorem{example}{Example}
\usepackage{mdframed} 
\usepackage{subcaption}
\usepackage[noend]{algorithmic}
\usepackage{algorithm}
\newtheorem*{remark}{Remark}
\newtheorem*{theorem}{Theorem}



\usepackage{titlesec}
% {left}{before-sep}{after-sep}
\titlespacing*{\section}{0pt}{1.2ex plus .2ex minus .2ex}{0.8ex plus .1ex}
\titlespacing*{\subsection}{0pt}{1.0ex plus .2ex minus .2ex}{0.5ex plus .1ex}

\usepackage{enumitem}
\setlist{nosep} % Removes all vertical space between list items and around lists




\usepackage[T1]{fontenc}
\usepackage{lmodern}

\usepackage{marginnote} 
\newcommand{\sid}[1]{%
    \marginnote{\tiny\raggedright\color{red}\textbf{SB:} #1}%
}


\title{Certification of True QBF Formulas in Expansion-Based Solving}
\author{
  Sidhant Bhavnani, Clemens Hofstadler, Martina Seidl \\
  \texttt{\{first\_name.last\_name\}@jku.at} \\
  Johannes Kepler University Linz, Austria \\
}

\date{}

\begin{document}

\maketitle
\begin{abstract}
    Expansion-based QBF solvers effectively decide false instances via universal expansion, yet the dual approach—deciding true instances via existential expansion—is less explored. This is because eliminating existential variables yields non-CNF formulas containing only universal variables, which standard SAT solvers cannot directly process. Consequently, certifying true results has remained an open problem. In this work, we present a proof format and checking algorithm for certifying true QBFs in the expansion-based setting, identifying and overcoming the core obstacles of non-clausal, universal-only expansions.
\end{abstract}


\section{Motivation}

Quantified Boolean formulas (QBFs) extend propositional logic with universal and
existential quantifiers, providing a natural encoding for problems in formal
verification, synthesis, and games~\cite{survey}. The QBF decision problem
is PSPACE-complete~\cite{survey}. Expansion-based solving is one of the two dominant
QBF solving paradigms, reducing the decision problem to propositional
satisfiability by eliminating quantifiers and delegating to a SAT solver.
Existing certification frameworks for expansion-based solving, such as
FERPModels~\cite{ferpmodels} and FERAT~\cite{ferat}, target only false QBFs via universal
expansion. Certifying true QBFs via existential expansion has remained an open
problem.

We present a proof format and checking algorithm for certifying true QBFs in the
expansion-based setting. Section~\ref{sec:preliminaries} establishes the necessary
background on expansion-based solving and existing certification frameworks.
Section~\ref{sec:existential} introduces existential expansion and the structural
difficulties it poses. Section~\ref{sec:challenge} identifies the two core
certification obstacles. Section~\ref{sec:format} presents the proof format and
the \textsc{CheckExpansion} algorithm.

\section{Preliminaries}
\label{sec:preliminaries}
A \emph{quantified Boolean formula} (QBF) in prenex conjunctive normal form (PCNF)
is a formula $\Phi = \Pi.\phi$, where the \emph{prefix} $\Pi = Q_1x_1\cdots Q_nx_n$
is a sequence of quantifier bindings with $Q_i \in \{\forall, \exists\}$, and the
\emph{matrix} $\phi$ is a propositional formula in CNF. We write $U_\Pi$ and $E_\Pi$
for the sets of universal and existential variables of $\Pi$ respectively, and
$<_\Pi$ for the quantifier order induced by $\Pi$.

A \emph{universal assignment} $\alpha$ is a total assignment to $U_\Pi$. The
\emph{instantiation} $\phi^\alpha$ is obtained by substituting each universal
variable with its value under $\alpha$, replacing each remaining existential
variable $e$ with a fresh \emph{annotated copy} $e^\tau$ where $\tau$ is the
restriction of $\alpha$ to $\{u \in U_\Pi \mid u <_\Pi e\}$, and simplifying. A
\emph{partial universal expansion} over a finite set $A$ of universal assignments
is the conjunction
\[
    \psi_\forall = \bigwedge_{\alpha \in A} \phi^\alpha.
\]
Since $\phi$ is in CNF, so is each $\phi^\alpha$, and hence so is $\psi_\forall$.
A QBF $\Phi$ is \emph{false} if and only if there exists a finite $A$ such that
$\psi_\forall$ is unsatisfiable~\cite{itjihad}.

\begin{example}
    \label{ex:running-false}
    Let $\Phi^\bot = \exists x \forall a \exists y.\, \phi^\bot$ where
    \[
        \phi^\bot = (x \vee a \vee y) \wedge (\neg x \vee \neg a \vee y) \wedge
        (\neg y).
    \]
    $\Phi^\bot$ is false. Expanding over $\alpha = \{a \leftarrow 0\}$ and $\alpha
        = \{a \leftarrow 1\}$ yields
    \[
        \psi_\forall = (x \vee y^0) \wedge (\neg y^0) \wedge (\neg x \vee y^1)
        \wedge (\neg y^1) = \bot
    \]
    where $y^0$ and $y^1$ are annotated copies of $y$ under $a \leftarrow 0$ and
    $a \leftarrow 1$ respectively.
\end{example}

When $\psi_\forall$ is passed to a SAT solver, the annotated variables lose their
connection to the original QBF variables. Expansion-based certification
frameworks~\cite{ferpmodels,ferat} restore this connection via two proof components:
\emph{expansion map} $\varepsilon$ mapping each proof variable to its annotated
QBF counterpart, encoded as \texttt{x} lines, and an \emph{origin map} recording
the originating matrix clause for each expansion clause, encoded as a single
\texttt{o} line. Figure~\ref{fig:false-proof} shows the QDIMACS file and the
expansion part of the proof file for $\Phi^\bot$.

\begin{figure}[t]
    \centering
    \begin{minipage}{0.45\columnwidth}
        \begin{mdframed}
            \begin{verbatim}
p cnf 3 3
e 1 0
a 2 0
e 3 0
1 2 3 0
-1 -2 3 0
-3 0

\end{verbatim}
        \end{mdframed}
        \subcaption{QDIMACS file for $\Phi^\bot$.}
    \end{minipage}
    \hfill
    \begin{minipage}{0.45\columnwidth}
        \begin{mdframed}
            \begin{verbatim}
x 1 0 1 0 0
x 3 0 3 0 -2 0
x 2 0 3 0 2 0
o 2 3 1 3 0
-1 2 0
-2 0
1 3 0
-3 0
\end{verbatim}
        \end{mdframed}
        \subcaption{Proof file for $\Phi^\bot$.}
    \end{minipage}
    \caption{QDIMACS and proof file for $\Phi^\bot$
        (Example~\ref{ex:running-false}).}
    \label{fig:false-proof}
\end{figure}

Variable \texttt{1} maps to $x$ with no annotation; variable \texttt{3} maps to
$y$ annotated by $a \leftarrow 0$; and variable \texttt{2} maps to $y$ annotated
by $a \leftarrow 1$. The \texttt{o} line records that the four expansion clauses
originate from matrix clauses \texttt{2}, \texttt{3}, \texttt{1}, and \texttt{3}
respectively. Given $\varepsilon$ and the origin map, the checker verifies
that each expansion clause is a legal instantiation of its originating matrix
clause under the recorded universal assignment, and delegates the refutation
checking of $\psi_\forall$ to a SAT solver and an external proof checker.

\section{Existential Expansion based Solving}
\label{sec:existential}
An \emph{existential assignment} $\sigma$ is a total assignment to $E_\Pi$. Given
$\sigma$, the \emph{instantiation} $\phi^\sigma$ is obtained by substituting each
existential variable with its value under $\sigma$, replacing each remaining
universal variable $a$ with a fresh \emph{annotated copy} $a^\tau$ where $\tau$ is
the restriction of $\sigma$ to $\{e \in E_\Pi \mid e <_\Pi a\}$, and simplifying.
A \emph{partial existential expansion} over a finite set $S$ of existential
assignments is the disjunction
\[
    \psi_\exists = \bigvee_{\sigma \in S} \phi^\sigma.
\]
A QBF $\Phi$ is \emph{true} if and only if there exists a finite $S$ such that
$\psi_\exists$ is valid~\cite{itjihad}, i.e.,
\[
    \forall \beta : U_\Pi \to \mathbb{B},\ \psi_\exists[\beta] = \top.
\]
Since validity of $\psi_\exists$ is equivalent to unsatisfiability of its negation,
\[
    \neg\psi_\exists = \bigwedge_{\sigma \in S} \neg\phi^\sigma,
\]
truth of $\Phi$ reduces to a SAT check on $\neg\psi_\exists$.


We now introduce the true QBF that serves as our running example for the remainder
of the paper.

\begin{example}
    \label{ex:running-true}
    Let $\Phi = \forall a_1 \exists e_2 \forall a_3 a_4 \exists e_5 e_6.\, \phi$
    where
    \begin{align*}
        \phi = \; & (\neg e_5 \vee \neg a_1 \vee a_4) \wedge
        (\neg e_5 \vee \neg e_2) \wedge
        (\neg e_5 \vee \neg a_3) \;\wedge                    \\
                  & (\neg e_6 \vee a_1) \wedge
        (\neg e_6 \vee e_2) \wedge
        (\neg e_6 \vee \neg a_3) \wedge
        (e_5 \vee e_6 \vee a_3).
    \end{align*}
    $\Phi$ is true.
\end{example}


The solver constructs a set $S$ of existential assignments and computes the
corresponding instantiations. For $\Phi$, the four non-trivial branches are:
\begin{align*}
    \phi^{000} & = (a_3^0) \qquad & \phi^{010} & = (\neg a_1 \vee a_4^0) \wedge (\neg a_3^0) \\
    \phi^{100} & = (a_3^1) \qquad & \phi^{101} & = (a_1) \wedge (\neg a_3^1)
\end{align*}
where the superscript on each annotated variable records the restriction of $\sigma$
to the existential variables to its left in $\Pi$. The partial existential expansion
is the disjunction
\[
    \psi_\exists = \phi^{000} \vee \phi^{010} \vee \phi^{100} \vee \phi^{101}.
\]

To check validity of $\psi_\exists$ via a SAT solver, we negate it to obtain
\[
    \neg\psi_\exists = \neg\phi^{000} \wedge \neg\phi^{010} \wedge \neg\phi^{100}
    \wedge \neg\phi^{101}.
\]
Since each $\phi^\sigma$ is in CNF, each $\neg\phi^\sigma$ is in DNF. For example,
\[
    \neg\phi^{010} = \neg(\neg a_1 \vee a_4^0) \vee \neg(\neg a_3^0)
    = (a_1 \wedge \neg a_4^0) \vee a_3^0,
\]
which is not clausal. Thus $\neg\psi_\exists$ is a conjunction of DNF formulas and  cannot be passed directly to a SAT solver.

To obtain an equisatisfiable CNF without exponential blowup, the solver applies the
Plaisted--Greenbaum (PG) transformation~\cite{pg} to each negated branch $\neg\phi^\sigma$.
For each clause $C_i \in \phi^\sigma$ with more than one literal, a fresh
\emph{auxiliary variable} $t_i$ is introduced satisfying $t_i \Rightarrow \neg C_i$,
contributing \emph{helper clauses}
\[
    H_i^\sigma = \{(t_i \vee \neg l) \mid l \in C_i\}.
\]
The negated branch $\neg\phi^\sigma$ is then encoded as a single \emph{NAND clause}
\[
    N^\sigma = (\neg\langle C_1 \rangle \vee \cdots \vee \neg\langle C_m \rangle)
\]
where $\langle C_i \rangle = t_i$ if $|C_i| > 1$ and $\langle C_i \rangle = l$ for
the unique literal $l \in C_i$ otherwise. The CNF formula passed to the SAT solver
is then
\[
    \psi_F = \bigwedge_{\sigma \in S} \left( N^\sigma \wedge \bigwedge_{|C_i|>1} H_i^\sigma \right).
\]
The auxiliary variables $t_i$ have no counterpart in $\phi$, and the NAND clause
$N^\sigma$ maps to multiple matrix clauses rather than one. This is the structure
the checker must accommodate.

Applying the PG transformation to $\Phi$, the negated branch $\neg\phi^{010}$
contains two clauses, $(\neg a_1 \vee a_4^0)$ and $(\neg a_3^0)$. The first has
more than one literal, so a fresh auxiliary variable $t_1$ is introduced, yielding
the helper clauses
\[
    H_1^{010} = (t_1 \vee a_1) \wedge (t_1 \vee \neg a_4^0)
\]
and the NAND clause
\[
    N^{010} = (\neg t_1 \vee a_3^0).
\]
The remaining branches $\phi^{000}$, $\phi^{100}$, and $\phi^{101}$ contain only
single-literal clauses, so no auxiliary variables are introduced and their NAND
clauses are
\[
    N^{000} = (\neg a_3^0), \quad N^{100} = (\neg a_3^1), \quad
    N^{101} = (\neg a_1 \vee a_3^1).
\]
The CNF formula passed to the SAT solver is therefore
\[
    \psi_F = N^{000} \wedge (N^{010} \wedge H_1^{010}) \wedge N^{100} \wedge N^{101}.
\]


\section{The Certification Challenge}
\label{sec:challenge}

Certifying the result of an expansion-based solver requires two things:
\begin{itemize}
    \item \emph{Correspondence}: a mapping between the clauses of the propositional
          formula passed to the SAT solver and the matrix clauses of $\phi$.
    \item \emph{Completeness}: a completeness check verifying that every matrix
          clause $C_i \in \phi$ is either accounted for by the expansion or
          satisfies $C_i^\sigma = \top$.
\end{itemize}
In the false case, correspondence is established clause by clause: the origin map
identifies which matrix clause each expansion clause was derived from, and
$\varepsilon$ maps each proof variable back to its annotated QBF counterpart,
allowing the checker to verify that each expansion clause is a legal instantiation
of its originating matrix clause under the recorded universal assignment.
Completeness is not a concern: it suffices to exhibit a subset of the universal
expansion that is unsatisfiable, as adding further clauses cannot change
unsatisfiability. In the true case, neither hold.

\begin{remark}[Broken Correspondence]
    \label{rem:correspondence}
    The PG transformation breaks the one-to-one correspondence between proof
    clauses and matrix clauses. Each NAND clause $N^\sigma$ maps to multiple
    matrix clauses, and the auxiliary variables $t_i$ have no counterpart in
    $\phi$. The checker cannot establish correspondence clause by clause as in
    the false case.
\end{remark}

\begin{remark}[Annotation Incompleteness]
    \label{rem:annotation}
    In the universal case, universal reduction guarantees the innermost quantifier
    block is always existential, so every existential variable has at least one
    universal variable to its left in $\Pi$ and receives a non-empty annotation.
    The full universal assignment $\alpha$ is therefore recoverable locally from
    the annotations on any single clause. In the true case, when the innermost
    quantifier block is existential, those existential variables annotate nothing
    --- their annotated copies carry an empty annotation $e^{\emptyset}$,
    indistinguishable from one another. The expansion map $\varepsilon$ alone does
    not determine $\sigma$, and the completeness check cannot be discharged.
\end{remark}

The two obstacles are interrelated. Broken correspondence means the checker cannot
determine which matrix clauses are accounted for by the expansion. Annotation
incompleteness means the checker cannot recover the full existential assignment
$\sigma$ needed to verify that the remaining clauses satisfy $C_i^\sigma = \top$.
Together, they prevent the completeness check from being discharged without
additional information supplied explicitly in the proof.

Consider the branch $\sigma = \{e_2 \leftarrow 0, e_5 \leftarrow 1, e_6
    \leftarrow 0\}$ for our running example $\Phi$, as derived in
Section~\ref{sec:existential}. The NAND clause $N^\sigma = (\neg t_1 \vee a_3^0)$
accounts for matrix clauses $C_1 = (\neg e_5 \vee \neg a_1 \vee a_4)$ and $C_7
    = (e_5 \vee e_6 \vee a_3)$. The remaining matrix clauses $C_2, C_3, C_4, C_5,
    C_6$ are not accounted for and must therefore satisfy $C_i^\sigma = \top$. For
example, $C_4 = (\neg e_6 \vee a_1)$ is eliminated because $e_6 \leftarrow 0
    \in \sigma$, satisfying $\neg e_6$. The proof format must therefore supply
additional information to bridge this gap.


\section{Proof Format and \textsc{CheckExpansion}}
\label{sec:format}

The proof must provide enough information to address both obstacles identified in
Section~\ref{sec:challenge}. Figure~\ref{fig:proof} shows the QDIMACS file and
the proof file for our running example $\Phi$. The proof consists of three line
types:
\begin{itemize}
    \item \texttt{x} lines encode the expansion map $\varepsilon$, mapping each
          proof variable to its annotated QBF counterpart. For example, \texttt{x 2 3
              0 3 4 0 2 0} records that proof variables \texttt{2} and \texttt{3} map to
          QBF variables \texttt{3} and \texttt{4} respectively, annotated by $e_2
              \leftarrow 1$.
    \item \texttt{n} lines encode NAND clause proof lines of the form
          \[
              \texttt{n} \;
              \underbrace{l_1 \cdots l_k}_{N^\sigma} \; 0 \;
              \underbrace{i_1 \cdots i_k}_{\eta_N} \; 0 \;
              \underbrace{\omega^\sigma}_{\text{witness}} \; 0
          \]
          where $l_1, \ldots, l_k$ are the literals of $N^\sigma$, $i_1, \ldots,
              i_k$ are the indices of their originating matrix clauses, and $\omega^\sigma$
          is the complete existential witness for branch $\sigma$. For example,
          \texttt{n 7 -3 0 1 7 0 -2 5 -6 0} encodes the NAND clause for branch
          $\sigma = \{e_2 \leftarrow 0, e_5 \leftarrow 1, e_6 \leftarrow 0\}$:
          literals \texttt{7} and \texttt{-3} are mapped to matrix clauses \texttt{1}
          and \texttt{7} respectively, and \texttt{-2 5 -6} is the witness
          $\omega^\sigma$.
    \item \texttt{h} lines encode helper clauses $H_i^\sigma$, one per multi-literal
          clause $C_i \in \phi^\sigma$, encoding the auxiliary variable definitions
          from the PG transformation. For example, \texttt{h 7 1 0} and \texttt{h 7
              -4 0} encode $H_1^\sigma$, the definition of auxiliary variable $t_1$
          introduced for clause $C_1 = (\neg e_5 \vee \neg a_1 \vee a_4)$.
\end{itemize}

\begin{figure}[t]
    \centering
    \begin{minipage}{0.45\columnwidth}
        \begin{mdframed}
            \begin{verbatim}
p cnf 6 7
a 1 0
e 2 0
a 3 4 0
e 5 6 0
-5 -1 4 0
-5 -2 0
-5 -3 0
-6 1 0
-6 2 0
-6 -3 0
5 6 3 0
\end{verbatim}
        \end{mdframed}
        \subcaption{QDIMACS file for $\Phi$.}
    \end{minipage}
    \hfill
    \begin{minipage}{0.45\columnwidth}
        \begin{mdframed}
            \begin{verbatim}
x 1 0 1 0 0
x 2 3 0 3 4 0 2 0
x 4 5 0 3 4 0 -2 0
n -3 0 7 0 -2 -5 -6 0
n 7 -3 0 1 7 0 -2 5 -6 0
h 7 1 0
h 7 -4 0
n -3 0 7 0 2 -5 -6 0
n -1 3 0 4 7 0 2 -5 6 0
\end{verbatim}
        \end{mdframed}
        \subcaption{Proof file for $\Phi$.}
    \end{minipage}
    \caption{QDIMACS and proof file for the running example $\Phi$
        (Example~\ref{ex:running-true}).}
    \label{fig:proof}
\end{figure}

\subsection{\textsc{CheckExpansion}}

The function $\mathrm{annotation}(v^\tau) = \tau$ extracts the annotation $\tau$
from an annotated variable $v^\tau$, treating it as a set of signed literals over
$E_\Pi$. We extend $\varepsilon$ to sets of literals pointwise: $\varepsilon(L) =
    \{\varepsilon(l) \mid l \in L\}$.

Given the proof, \textsc{CheckExpansion} iterates over every NAND clause $N^\sigma$
and performs three steps:
\begin{enumerate}
    \item \emph{Check Literals}: for each literal $l \in N^\sigma$, retrieve the
          originating matrix clause $C_i = \phi[\eta_N(l)]$. Here $\mathrm{universal}
              (C_i)$ denotes the set of universal literals of $C_i$ as they appear in
          $\phi$. If $l$ is an auxiliary variable $t_i$, first verify that the helper
          clauses $H_i^\sigma$ correctly encode $t_i \Rightarrow \neg C_i$ under
          $\varepsilon$, then collect the non-auxiliary literals of $H_i^\sigma$;
          otherwise use $l$ directly. Rename the collected literals via $\varepsilon$,
          stripping annotations to recover plain QBF variables, and verify they match
          $\mathrm{universal}(C_i)$. Accumulate the existential annotations into a
          partial witness $\omega_{\mathrm{partial}}$ via $\mathrm{annotation}(\cdot)$,
          where $\mathrm{annotation}(v^\tau) = \tau$ extracts $\tau$ as a set of signed
          literals over $E_\Pi$.
    \item \emph{Validate Witness}: extend $\omega_{\mathrm{partial}}$ with the
          supplied witness $\omega^\sigma$, verifying that $\omega^\sigma$ is
          consistent with $\omega_{\mathrm{partial}}$, that all variables in
          $\omega^\sigma$ are existential variables of $\Pi$, and that the result
          $\omega$ satisfies $\mathrm{dom}(\omega) = E_\Pi$.
    \item \emph{Check Elimination}: for every matrix clause $C_i \in \phi$ not
          referenced by $\eta_N$, verify that $C_i^{\omega} = \top$, i.e., at
          least one existential literal of $C_i$ is satisfied by $\omega$.
\end{enumerate}

\begin{algorithm}
    \caption{\textsc{CheckExpansion}}
    \begin{algorithmic}[1]
        \REQUIRE QBF $\Pi.\phi$, proof with $\varepsilon$, NAND clauses, helper
        clauses, $\eta_N$, witnesses $\omega^\sigma$
        \ENSURE \textit{true} if all NAND clauses are valid, \textit{false} otherwise

        \FOR{each NAND clause $N^\sigma$ in proof}

        \STATE $\omega_{\mathrm{partial}} \leftarrow \emptyset$

        \STATE $\triangleright$ \textit{Step 1 --- Check Literals}
        \FOR{each $l \in N^\sigma$}
        \STATE $C_i \leftarrow \phi[\eta_N(l)]$
        \IF{$\textsc{IsAuxiliary}(l)$}
        \STATE $\mathit{lits} \leftarrow \textsc{GetNonAuxiliaryLiterals}(H_i^\sigma)$
        \ELSE
        \STATE $\mathit{lits} \leftarrow \{l\}$
        \ENDIF
        \IF{$\varepsilon(\mathit{lits}) \neq \mathrm{universal}(C_i)$}
        \RETURN \textit{false}
        \ENDIF
        \STATE $\omega_{\mathrm{partial}} \leftarrow \omega_{\mathrm{partial}}
            \cup \{\mathrm{annotation}(k) \mid k \in \mathit{lits}\}$
        \ENDFOR

        \STATE $\triangleright$ \textit{Step 2 --- Validate Witness}
        \FOR{each $l \in \omega^\sigma$}
        \IF{$\mathrm{Var}(l) \notin E_\Pi$}
        \RETURN \textit{false}
        \ENDIF
        \IF{$\neg l \in \omega_{\mathrm{partial}}$}
        \RETURN \textit{false}
        \ENDIF
        \STATE $\omega_{\mathrm{partial}} \leftarrow \omega_{\mathrm{partial}} \cup \{l\}$
        \ENDFOR
        \IF{$\mathrm{dom}(\omega_{\mathrm{partial}}) \neq E_\Pi$}
        \RETURN \textit{false}
        \ENDIF
        \STATE $\omega \leftarrow \omega_{\mathrm{partial}}$

        \STATE $\triangleright$ \textit{Step 3 --- Check Elimination}
        \FOR{each $C_i \in \phi$ not referenced by $\eta_N$}
        \IF{$C_i^\omega \neq \top$}
        \RETURN \textit{false}
        \ENDIF
        \ENDFOR

        \ENDFOR
        \RETURN \textit{true}
    \end{algorithmic}
\end{algorithm}

The origin map $\eta_N$ and helper clauses $H_i^\sigma$ resolve broken
correspondence by restoring a clause-level mapping to $\phi$. The witness
$\omega^\sigma$ resolves annotation incompleteness by supplying the full
existential assignment required to discharge the completeness check.



\section{Conclusion}\label{sec:conclusion}
We have presented a proof format and checking algorithm for certifying true QBFs
in the expansion-based setting. The two core obstacles --- broken correspondence
introduced by the Plaisted--Greenbaum transformation, and annotation incompleteness
arising when the innermost quantifier block is existential --- are addressed by
augmenting the expansion proof with an origin map $\eta_N$, helper clauses
$H_i^\sigma$, and a per-branch existential witness $\omega^\sigma$. The
\textsc{CheckExpansion} algorithm verifies each NAND clause in three steps,
culminating in a completeness check that enforces the key invariant: every matrix
clause is either referenced by a NAND clause or eliminated by the witness.

Several directions remain for future work. The proof format and
\textsc{CheckExpansion} algorithm presented here provide the theoretical foundation
for a practical proof checker; implementing and evaluating this checker on
expansion-based QBF benchmarks is a natural next step. Beyond certification,
FERPModels already supports Herbrand function extraction for false QBFs; extending
this to Skolem function extraction for true QBFs from verified existential
expansion proofs would provide concrete witnesses to the truth of a QBF,
broadening the utility of expansion-based certification in synthesis and
verification applications.

\bibliographystyle{plain}
\bibliography{references}
\end{document}
